\section{Flow Matching, Rectified Flow \& Flux.1}\label{sec:flow}
\lab{Flow Matching \textperiodcentered\ straight paths \textperiodcentered\ velocity objective \textperiodcentered\ Rectified Flow \textperiodcentered\ Flux.1 12B \textperiodcentered\ OT-CFM}

\diagpair{%
\adjustbox{max width=\linewidth}{%
\begin{tikzpicture}[x=1mm,y=1mm,every node/.style={font=\scriptsize}]
\node[box,draw=mute] (x0) at (0,10) {$x_0 \sim \mathcal{N}(0, \mathbf{I})$};
\node[box,draw=acc] (x1) at (30,10) {$x_1 \sim q(x_1)$};
% Curved diffusion path
\draw[draw=acc2,thick,dashed] (x0) to[bend left=40] node[midway,above] {Curved Diffusion} (x1);
% Straight Flow Matching path
\draw[arr,acc3,very thick] (x0) -- (x1) node[midway,below] {Straight Path ($v_t$)};
\end{tikzpicture}%
}
}{%
\textbf{Straight Probability Paths.} Diffusion trajectories curve stochastically through latent space, requiring small steps.\\[2pt]
{\color{acc2}$\blacktriangleright$} \textbf{Constant Velocity Target}:
$x_t = (1 - t) x_0 + t x_1 \implies v_t = \frac{d x_t}{dt} = x_1 - x_0$\\[2pt]
{\color{acc2}$\blacktriangleright$} \textbf{Single-step Euler}: Straight paths make 1-to-4 step generation possible without distillation.
}

\begin{mem}[Flow Matching vs Diffusion Comparison]
\begin{tabularx}{\linewidth}{@{}p{18mm} X X@{}}
\toprule
\textbf{Property} & \textbf{Standard Diffusion (DDPM/EDM)} & \textbf{Flow Matching / Rectified Flow}\\
\midrule
\textbf{Trajectory} & Curved probability paths & \textbf{Straight} linear trajectories ($x_t = (1-t)x_0 + tx_1$)\\
\textbf{Prediction Target} & Noise $\epsilon_t$ or data $x_0$ & Constant velocity vector field $v_t = x_1 - x_0$\\
\textbf{Time Parameter} & Discrete $t \in [1, 1000]$ or continuous $\sigma$ & Normalized interval $t \in [0, 1]$\\
\textbf{Sampling Step} & Complex variance-adjusted ODE step & Simple linear update: $x_{t+\Delta t} = x_t + \Delta t \cdot v_\theta(x_t, t)$\\
\textbf{Few-Step Quality} & Blurry / degraded under 10 steps & Crisp, detailed generation in 4--8 steps\\
\bottomrule
\end{tabularx}
\end{mem}

\begin{qa}[Rectified Flow Mechanics \& Flux Architecture]
\Q{How does Rectified Flow straighten sampling paths, and what is Reflowing?}{
When training on random pairs of $(x_0, x_1)$ ($x_0$ random noise, $x_1$ data image), trajectories can cross each other in state space. Where paths cross, the vector field must predict the average velocity of diverging paths, inducing local curvature.
\textbf{Reflowing Procedure (2-Rectified Flow)}:
\begin{enumerate}[leftmargin=*,itemsep=0.8pt]
\item Generate synthetic pairs $(x_0, \hat{x}_1)$ by simulating the trained 1-rectified flow ODE from $t=0$ to $t=1$.
\item Train a new flow matching model on these deterministic coupled pairs $(x_0, \hat{x}_1)$.
\end{enumerate}
Because the pairs come from an ODE whose paths cannot cross (uniqueness of ODE solutions), the new coupling has less crossing and the learned paths are straighter; Liu et al.\ prove reflow never increases transport cost. In practice 2-rectified flow plus distillation gives good samples in \textbf{1 to 2 Euler steps}.}
\Q{What is Optimal Transport Flow Matching (OT-CFM)?}{
Tong et al.\ (2023) observe that standard flow matching pairs random Gaussian noise $x_0$ with random data $x_1$ independently, producing long, crossing trajectories.
\textbf{Optimal Transport CFM}: Couples pairs $(x_0, x_1)$ according to the optimal transport plan $\pi$ minimizing dynamic kinetic energy:
$$\mathcal{W}_2^2(p_0, p_1) = \inf_\pi \int \|x_0 - x_1\|_2^2 \, d\pi(x_0, x_1)$$
Solve exact OT (Hungarian / EMD) within each mini-batch and re-pair noise with data by that plan. Paths cross less, training variance drops and fewer ODE steps are needed; the benefit shrinks in high dimensions where mini-batch OT is a rough approximation.}
\Q{Walk through the 12B parameter architecture of Flux.1.}{
Flux.1 (Black Forest Labs) is an MM-DiT in the style of SD3, with no cross-attention:
\begin{enumerate}[leftmargin=*,itemsep=0.8pt]
\item \textbf{Double Stream Blocks (19 blocks)}: Image and text tokens keep \emph{separate weights} (norms, QKV, MLP) but run \emph{one joint attention} over the concatenated sequence, so each modality reads the other while keeping specialized parameters.
\item \textbf{Single Stream Blocks (38 blocks)}: Concatenates text and visual tokens into a single unified residual stream, running bidirectional self-attention across all tokens simultaneously.
\item \textbf{Dual Text Conditioning}: \textbf{T5-XXL} token sequence (prompt structure, spelling) enters the text stream; the pooled \textbf{CLIP ViT-L} vector is added to the timestep embedding for adaLN modulation (global style).
\item \textbf{Rotary Positional Embeddings (2D RoPE)}: Replaces fixed absolute coordinates, enabling training and inference across any aspect ratio and resolution up to 2.0 megapixels without tiling.
\end{enumerate}}
\end{qa}

\begin{trap}[Flow Matching Gotchas]
\begin{itemize}[leftmargin=*,itemsep=0.8pt]
\item \textbf{Direction of $t$}: In diffusion literature, $t=0$ is clean data and $t=T$ is pure noise. In Flow Matching literature, conventions are often \textbf{inverted}: $t=0$ is pure Gaussian noise ($x_0 \sim \mathcal{N}(0, \mathbf{I})$) and $t=1$ is target data ($x_1$). Check the boundary conditions of your solver ($0 \to 1$ vs $1 \to 0$).
\item \textbf{CFG formulation in Flow Matching}: Naive CFG extrapolation on velocity fields can overshoot data manifolds. Flux uses guidance distillation or interval-constrained CFG to prevent color explosion.
\end{itemize}
\end{trap}
